\section*{Erd\H{o}s Problem 698: a growing common divisor of binomial coefficients}

\paragraph{Statement and claim.}
For \(2\leq i<j\leq n/2\), is
\(\gcd(\binom ni,\binom nj)\) uniformly bounded below by a
function tending to infinity with \(n\)?
Yes: one may take \(h(n)=\sqrt n/6\).
We reconstruct Bergman's argument, including its stronger bound
\[
 \gcd\left(\binom ni,\binom nj\right)
 \geq\sqrt n\,\frac{2^{i-7/2}}{i\sqrt{i-1}}.
\]
Source: George M. Bergman, \emph{On common divisors of multinomial
coefficients}, Section 2, Theorem 2,
\url{https://math.berkeley.edu/~gbergman/papers/nomial.pdf}.
See \url{https://www.erdosproblems.com/698}.

\paragraph{Common multiples with controlled variation.}
Write \(A=\binom ni\), \(B=\binom nj\), and \(L=\operatorname{lcm}(A,B)\).
For \(0\leq h\leq i\), define
\[
 Q_h=A\binom ih\binom{n-i}{j-i+h}
     =B\binom j{i-h}\binom{n-j}h.
\]
The equality follows by expanding factorials; equivalently, it
counts ordered pairs of subsets of sizes \(i,j\) with exactly
\(h\) elements of the first outside the second.
Both expressions are integer multiples of \(A,B\), so \(L\mid Q_h\).
Since \(i\geq2\), all of \(Q_0,Q_1,Q_2\) are available.
It follows that \(L^2\) divides the integer
\[
 D=(i-1)Q_1^2-2iQ_0Q_2.
\]
An explicit cancellation gives
\[
 D=B^2\binom j{i-2}^{\!2}
       \frac{(j-i+2)(n-j)(n-i+1)}{i-1}>0.                      \tag{1}
\]
For clarity, after factoring
\(B^2\binom j{i-2}^2(j-i+2)(n-j)/(i-1)\),
the remaining expression is
\[
 (j-i+2)(n-j)-(j-i+1)(n-j-1)=n-i+1,
\]
which proves (1). Positivity is crucial: divisibility now yields
\(L^2\leq D\).

\paragraph{Estimating the positive common multiple.}
Since \(j\leq n/2\), for every \(0\leq r<i-2\),
\((j-r)/(n-r)\leq1/2\). Thus
\[
 \binom j{i-2}\leq2^{-(i-2)}\binom n{i-2}.
\]
Also \(j-i+2\leq n/2\), \(n-j<n\), and \(n-i+1<n\).
Substitution in (1) gives
\[
 L^2<
 \frac{n^3}{2^{2i-3}(i-1)}
 B^2\binom n{i-2}^{\!2}.
\]
Using \(\gcd(A,B)=AB/L\), we obtain
\[
 \gcd(A,B)>
 \frac{(n-i+2)(n-i+1)}{i(i-1)}
 \sqrt{\frac{2^{2i-3}(i-1)}{n^3}}
 \geq\sqrt n\,\frac{2^{i-7/2}}{i\sqrt{i-1}}.
\]
The last step uses \(n-i+1,n-i+2\geq n/2\).

\paragraph{A bound independent of \(i\).}
For \(i=2\) the coefficient is \(1/(4\sqrt2)>1/6\);
for \(i=3\) it is \(1/6\).
For \(i\geq3\), the ratio of the coefficient at \(i+1\) to that
at \(i\) is
\[
 \frac{2\sqrt{i(i-1)}}{i+1}>1,
\]
because \(4i(i-1)>(i+1)^2\).
The minimum is therefore \(1/6\), proving the stated choice of
\(h(n)\). For the small \(n\) admitting no pair of indices, the
assertion is vacuous.

\paragraph{Scope.}
This is a full elementary reconstruction of the established
positive solution; it uses no asymptotic number-theoretic theorem.

\paragraph{Completion estimate.}
\textbf{COMPLETION: 100\%}
